\chapter*{Abstract}
\addcontentsline{toc}{chapter}{Abstract}

Large language models generate fluent text even when their answers are
semantically inconsistent or factually unsupported, and they rarely
signal when an answer should not be trusted. This thesis studies
\emph{semantic entropy} as a training-free method for hallucination
detection and selective abstention in short-answer question answering.

Semantic entropy measures uncertainty over meanings rather than over
surface strings. Sampled answers are grouped into semantic equivalence
classes using bidirectional Natural Language Inference entailment, and
Shannon entropy is computed over cluster frequencies. The method is
designed to detect \emph{confabulations}---answers that are both
incorrect and semantically inconsistent across alternative generations
from the same prompt---while making no claim about high-confidence
systematic errors that fall outside this failure mode.

We implement discrete semantic entropy on two structurally distinct
tasks: TriviaQA, an open-domain factual question-answering benchmark,
and SVAMP, a short-answer arithmetic reasoning dataset. Experiments are
organized in two phases on NHR@FAU high-performance computing resources.
Phase~1 uses Mistral-7B with a DeBERTa NLI backend to validate the
end-to-end pipeline at modest scale. Phase~2 extends the evaluation to
Llama-3.1-70B and stronger semantic judges on A100 and H100 hardware.
Uncertainty quality is measured with the area under the receiver
operating characteristic curve (AUROC), the area under the
rejection-accuracy curve (AURAC), and rejection-accuracy curves that
compare discrete semantic entropy against surface-form entropy baselines.

% TODO: replace the paragraph below with empirical findings once
% Phase 1 and Phase 2 HPC results are available.
Preliminary infrastructure validation confirms that the pipeline
produces schema-valid output records across both datasets and that
correctness labeling and entropy scoring are end-to-end reproducible.
Full empirical results comparing semantic entropy against surface-form
baselines, together with a manual audit of semantic clustering behavior,
are reported in Chapter~6 following completion of the HPC runs.

We additionally scope two Phase~2 extensions---Semantic Entropy Probes,
which approximate semantic entropy from hidden model states, and Kernel
Language Entropy, which replaces hard clustering with a kernel-based
entropy objective---as cost-versus-quality comparators once the discrete
baseline is established. The implementation, JSONL output schema, and
evaluation scripts are made available as reproducible artefacts for
future work on semantic uncertainty in large language models.
